\chapter{Local additions and manifolds of mappings}
\label{chap:mapping}

Glöckner and Neeb build the manifold structure of $C^\ell(M, N)$, for $M$ compact and $N$ with a
local addition, in \cite[\S3.8, \S4.1, \S4.2]{GN26}. This chapter follows them: local additions
and the ways to produce them (open subsets, products, normalization, ambient formulas on
submanifolds, real analytic local additions), the manifolds $C(K, N)$ and $C^\ell(M, N)$ with
their exponential law, evaluation maps and tangent spaces, and the compact-open $C^\ell$ topology,
described both through iterated tangent maps \cite[Def.~4.1.2]{GN26} and through jets on compact
sets in charts. The models of $C^\ell(M, N)$ vary with the point, so these manifolds are atlases
with varying models. As in the rest of the library, \emph{no completeness} of the models is
assumed.

\section{Local additions}

\begin{definition}[Local additions {\cite[Def.~3.8.11]{GN26}}]
  \label{def:local-addition}
  \lean{LocallyConvex.Manifolds.LocalAddition, LocallyConvex.Manifolds.LocalAddition.toFun,
    LocallyConvex.Manifolds.LocalAddition.source,
    LocallyConvex.Manifolds.LocalAddition.smoothMBMOn_toFun,
    LocallyConvex.Manifolds.LocalAddition.toOpenPartialHomeomorph_apply,
    LocallyConvex.Manifolds.LocalAddition.IsNormalized,
    LocallyConvex.Manifolds.TangentBundle.vec}
  \leanok
  \uses{def:tangent-bundle, def:contmdiff, prop:man-product}
  A \emph{local addition} on $M$ is a smooth map $\Sigma \colon \Omega \to M$ on an open subset
  $\Omega \subseteq \T M$ containing the zero section, with $\Sigma(0_x) = x$, such that
  $\theta = (\pi, \Sigma) \colon \Omega \to M \times M$ is a smooth diffeomorphism onto an open
  subset of $M \times M$. It is \emph{normalized} if $T_{0_x}(\Sigma|_{\T_x M}) = \mathrm{id}$
  for all $x$; in the preferred chart $\varphi_x$ at $x$, this reads
  $\dd(\varphi_x \circ \Sigma(x, \cdot))(0)(w) = w$. In Lean the data is the map $\theta$, an
  open partial homeomorphism $\T M \to M \times M$ with source $\Omega$, smooth with smooth
  inverse, whose first component is the bundle projection.
\end{definition}

\begin{proposition}[Open subsets of the model space]
  \label{prop:open-local-addition}
  \lean{LocallyConvex.Manifolds.opens_atlas_apply,
    LocallyConvex.Manifolds.tangentCoordChange_opens,
    LocallyConvex.Manifolds.chartAt_tangentBundle_opens_apply,
    LocallyConvex.Manifolds.addOpen, LocallyConvex.Manifolds.smoothMBMOn_addOpen,
    LocallyConvex.Manifolds.smoothMBMOn_addOpen_symm,
    LocallyConvex.Manifolds.openLocalAddition,
    LocallyConvex.Manifolds.isNormalized_openLocalAddition}
  \leanok
  \uses{def:local-addition}
  Let $U$ be an open subset of $E$. Then $\Sigma(x, v) = x + v$ on
  $\Omega = \{(x, v) \in \T U : x + v \in U\}$ is a normalized local addition on $U$, and
  $\theta(x, v) = (x, x + v)$ maps $\Omega$ onto $U \times U$.
\end{proposition}

\begin{proof}
  \leanok
  The charts of $U$ are the inclusion, so $\T U = U \times E$ with the identity as coordinate
  changes. In these coordinates $\theta$ and $\theta^{-1}(x, y) = (x, y - x)$ are restrictions
  of continuous linear maps.
\end{proof}

\begin{proposition}[Normalization {\cite[Remark~3.8.16]{GN26}}]
  \label{prop:normalize-local-addition}
  \lean{LocallyConvex.Manifolds.LocalAddition.alpha, LocallyConvex.Manifolds.LocalAddition.beta,
    LocallyConvex.Manifolds.LocalAddition.alpha_beta,
    LocallyConvex.Manifolds.LocalAddition.beta_alpha,
    LocallyConvex.Manifolds.smoothMBM_of_coordAt,
    LocallyConvex.Manifolds.LocalAddition.smoothMBM_normMap,
    LocallyConvex.Manifolds.LocalAddition.smoothMBM_normMapInv,
    LocallyConvex.Manifolds.LocalAddition.normEquiv,
    LocallyConvex.Manifolds.LocalAddition.normalize,
    LocallyConvex.Manifolds.LocalAddition.normalize_toFun,
    LocallyConvex.Manifolds.LocalAddition.isNormalized_normalize}
  \leanok
  \uses{def:local-addition}
  Let $\Sigma$ be a local addition on $M$. Then $\alpha_x = T_{0_x}(\Sigma|_{\T_x M})$ is
  invertible for every $x$, the map $F \colon \T M \to \T M$, $F(v) = \alpha_{\pi v}^{-1} v$, is a
  diffeomorphism, and $\Sigma \circ F$ is a normalized local addition.
\end{proposition}

\begin{proof}
  \leanok
  In the preferred chart $\varphi$ at $x$, the map $v \mapsto \varphi(\Sigma(x, v))$ has the local
  inverse $z \mapsto \theta^{-1}(x, \varphi^{-1} z)$ read in the tangent chart, so $\alpha_x$ is
  invertible by the chain rule. In an arbitrary chart, $\alpha_x$ and $\alpha_x^{-1}$ are
  conjugated by the derivative of the transition map and are partial derivatives of the chart
  representatives of $\theta$ and $\theta^{-1}$; hence $F$ and $F^{-1}$ are smooth. Glöckner and
  Neeb argue instead that $\theta$ restricts to diffeomorphisms of the fibres; we do not claim
  this argument is new.
\end{proof}

\begin{lemma}[The tangent bundle of a product]
  \label{lem:tangent-prod}
  \lean{LocallyConvex.Manifolds.prodChart_mem_atlas,
    LocallyConvex.Manifolds.tangentCoordChange_prod, LocallyConvex.Manifolds.tangentProdMk,
    LocallyConvex.Manifolds.tangentProdSplit, LocallyConvex.Manifolds.tangentProdSplit_eq,
    LocallyConvex.Manifolds.smoothMBM_tangentProdSplit,
    LocallyConvex.Manifolds.chartAt_tangentProdMk,
    LocallyConvex.Manifolds.continuous_tangentProdMk,
    LocallyConvex.Manifolds.smoothMBM_tangentProdMk,
    LocallyConvex.Manifolds.tangentProdHomeomorph}
  \leanok
  \uses{thm:man-tangent-manifold, thm:man-tangent-map-smooth, prop:man-product}
  The map $\T(M_1 \times M_2) \to \T M_1 \times \T M_2$, $v \mapsto (\T \mathrm{pr}_1(v),
  \T \mathrm{pr}_2(v))$, is a smooth diffeomorphism.
\end{lemma}

\begin{proof}
  \leanok
  The coordinate changes of $\T(M_1 \times M_2)$ between product charts are the products of the
  coordinate changes of $\T M_1$ and $\T M_2$. Hence, in charts, the inverse map
  $(v_1, v_2) \mapsto (v_1, v_2)$ is the regrouping $((a, b), (c, d)) \mapsto ((a, c), (b, d))$
  of coordinates, which is continuous linear; the map itself is smooth by
  Theorem~\ref{thm:man-tangent-map-smooth}.
\end{proof}

\begin{proposition}[Products of local additions {\cite[Lemma~3.8.17]{GN26}}]
  \label{prop:prod-local-addition}
  \lean{LocallyConvex.Manifolds.SmoothMBMOn.prodMap,
    LocallyConvex.Manifolds.smoothMBM_prodProdProdComm,
    LocallyConvex.Manifolds.LocalAddition.prodTheta, LocallyConvex.Manifolds.LocalAddition.prod}
  \leanok
  \uses{def:local-addition, lem:tangent-prod}
  If $\Sigma_1$ and $\Sigma_2$ are local additions on $M_1$ and $M_2$, then, identifying
  $\T(M_1 \times M_2)$ with $\T M_1 \times \T M_2$, $\Sigma_1 \times \Sigma_2$ is a local addition on
  $M_1 \times M_2$.
\end{proposition}

\begin{proof}
  \leanok
  The map $\theta$ is $\theta_1 \times \theta_2$ composed with the identification of
  Lemma~\ref{lem:tangent-prod} and the regrouping
  $(M_1 \times M_1) \times (M_2 \times M_2) \cong (M_1 \times M_2) \times (M_1 \times M_2)$, all of
  which are smooth with smooth inverses.
\end{proof}

\begin{proposition}[Local additions on embedded submanifolds from ambient formulas]
  \label{prop:ambient-local-addition}
  \lean{LocallyConvex.Manifolds.IsEmbeddedSubmanifold.contMDiffMBOn_of_val,
    LocallyConvex.Manifolds.IsEmbeddedSubmanifold.ambVec,
    LocallyConvex.Manifolds.IsEmbeddedSubmanifold.ofAmb,
    LocallyConvex.Manifolds.IsEmbeddedSubmanifold.coe_coordAt_slice,
    LocallyConvex.Manifolds.IsEmbeddedSubmanifold.smoothMBM_ambient,
    LocallyConvex.Manifolds.IsEmbeddedSubmanifold.contMDiffMBOn_of_ambient,
    LocallyConvex.Manifolds.AmbientAddition,
    LocallyConvex.Manifolds.AmbientAddition.toLocalAddition,
    LocallyConvex.Manifolds.AmbientAddition.toLocalAddition_toFun,
    LocallyConvex.Manifolds.AmbientAddition.isNormalized_toLocalAddition}
  \leanok
  \uses{def:local-addition, def:embedded-submanifold, def:tangent-space-submanifold}
  Let $X \subseteq F$ be an embedded submanifold modeled on a closed subspace $S$. Let
  $\sigma \colon \Omega \to F$ and $\tau \colon V \to F$ be smooth on open subsets of $F \times F$
  such that, for $p \in X$, $(p, 0) \in \Omega$, $\sigma(p, 0) = p$, and $w \mapsto \sigma(p, w)$ is a
  bijection from $\{w \in T_pX : (p, w) \in \Omega\}$ onto $\{q \in X : (p, q) \in V\}$ with inverse
  $q \mapsto \tau(p, q)$. Then $\Sigma(v) = \sigma(\pi v, v)$, with $v$ read in $F$, is a local
  addition on $X$, normalized if $\dd\sigma(p, \cdot)(0) = \mathrm{id}$ on $T_pX$.
\end{proposition}

\begin{proof}
  \leanok
  A map into $X$ is smooth when it is smooth into $F$: in a slice chart $\varphi$ its local
  representative is the corestriction to $S$ of $\varphi$ composed with the map into $F$. In the
  same chart, the fiber coordinate of a tangent vector $v$ at $p$ is $\dd\varphi(p) w$, where
  $w \in T_pX \subseteq F$ is its ambient vector; so the ambient map $\T X \to F \times F$,
  $v \mapsto (p, w)$, is smooth, and a map into $\T X$ is smooth when its composition with the
  ambient map is. Then $\theta = (\pi, \Sigma)$ and $\theta^{-1}(p, q) = (p, \tau(p, q))$ are
  smooth. Glöckner and Neeb obtain local additions from Lie group structures and from sprays
  \cite[Prop.~3.8.12, Prop.~3.8.13]{GN26}; we do not claim this argument is new.
\end{proof}

\begin{proposition}[Injectivity near a compact set {\cite[Prop.~3.8.18, Cor.~3.8.19]{GN26}}]
  \label{prop:local-diffeo-compact}
  \lean{LocallyConvex.Manifolds.exists_isOpen_injOn_of_isCompact,
    LocallyConvex.Manifolds.IsLocalHomeomorph.exists_openPartialHomeomorph_of_isCompact,
    LocallyConvex.Manifolds.ContMDiffMBOn.of_local,
    LocallyConvex.Manifolds.OpenPartialHomeomorph.contMDiffMBOn_of_local,
    LocallyConvex.Manifolds.OpenPartialHomeomorph.contMDiffMBOn_symm_of_local,
    LocallyConvex.Manifolds.exists_contMDiffMB_diffeo_of_isCompact,
    LocallyConvex.Manifolds.exists_smoothMBM_diffeo_of_isCompact}
  \leanok
  \uses{def:contmdiff}
  A continuous map $f \colon X \to Y$ into a Hausdorff space, injective on a compact set $K$ and
  locally injective at the points of $K$, is injective on an open neighborhood of $K$. Hence a
  local $C^k$-diffeomorphism ($k \in \N$ or $k = \infty$) into a Hausdorff manifold that is
  injective on a compact set $K$ is a $C^k$-diffeomorphism from an open neighborhood of $K$ onto
  an open set.
\end{proposition}

\begin{proof}
  \leanok
  The set $\{(x, y) : f(x) = f(y) \Rightarrow x = y\}$ is a neighborhood of each point of
  $K \times K$: of $(x, x)$ by local injectivity, of $(x, y)$ with $x \ne y$ because
  $f(x) \ne f(y)$ have disjoint neighborhoods. By the Wallace lemma it contains $A \times B$ for
  neighborhoods $A$, $B$ of $K$, and $f$ is injective on the interior of $A \cap B$. Glöckner and
  Neeb argue with nets; we do not claim this argument is new. For the second statement, the
  restriction is open and continuous, hence a homeomorphism onto its open image, and its inverse
  agrees near each point with a local inverse of $f$. The real analytic case is
  Corollary~\ref{cor:local-diffeo-analytic}.
\end{proof}

\begin{proposition}[Injectivity near a $\sigma$-compact set {\cite[Prop.~3.8.20, Cor.~3.8.21]{GN26}}]
  \label{prop:local-diffeo-sigma-compact}
  \lean{LocallyConvex.Manifolds.IsSigmaCompact.exists_compact_exhaustion,
    LocallyConvex.Manifolds.exists_isOpen_injOn_of_isSigmaCompact_image,
    LocallyConvex.Manifolds.exists_isOpen_injOn_of_isSigmaCompact,
    LocallyConvex.Manifolds.exists_openPartialHomeomorph_of_injOn,
    LocallyConvex.Manifolds.exists_openPartialHomeomorph_of_isSigmaCompact,
    LocallyConvex.Manifolds.locallyCompactSpace_of_finiteDimensional,
    LocallyConvex.Manifolds.exists_contMDiffMB_diffeo_of_isSigmaCompact,
    LocallyConvex.Manifolds.exists_smoothMBM_diffeo_of_isSigmaCompact}
  \leanok
  \uses{prop:local-diffeo-compact}
  Let $f \colon X \to Y$ be continuous, $Y$ Hausdorff with every point having a compact
  neighborhood, and $A \subseteq X$ $\sigma$-compact such that $f(A)$ is closed, $f|_A$ is an
  embedding and $f$ is locally injective at the points of $A$. Then $f$ is injective on an open
  neighborhood of $A$. Hence a local $C^k$-diffeomorphism ($k \in \N$ or $k = \infty$) into a
  Hausdorff, locally compact manifold (for instance a finite-dimensional one) which restricts to
  such an embedding on $A$ is a $C^k$-diffeomorphism from an open neighborhood of $A$ onto an open
  set.
\end{proposition}

\begin{proof}
  \leanok
  Take a compact exhaustion $\emptyset = C_0 \subseteq C_1 \subseteq \cdots$ of a neighborhood of
  $f(A)$ with $C_n \subseteq \operatorname{int} C_{n+1}$. The ``annuli''
  $O_j = \operatorname{int} C_{j+1} \setminus C_{j-1}$ meet only for $|i - j| \le 1$. Let $V_j$
  be an open neighborhood of the compact set $A \cap f^{-1}(C_{j+1})$ on which $f$ is injective
  (Proposition~\ref{prop:local-diffeo-compact}); then $f$ is injective on
  $U = \bigcup_j f^{-1}(O_j) \cap V_j \cap V_{j+1}$. Glöckner and Neeb use their Lemma~A.5.5
  instead of the annuli, and the Hausdorff property of the domain, which is not needed here. The
  statement of \cite[Cor.~3.8.21]{GN26} concludes $K \subseteq U$, a misprint for $A \subseteq U$;
  its finite-dimensionality hypothesis is only used through local compactness of the target. We
  do not claim this argument is new.
\end{proof}

\section{Real analytic local additions}

Real analyticity is meant in Milnor's sense (Chapter~\ref{chap:analytic}): local restrictions of
complex analytic maps on complexifications, with no completeness of the model.

\begin{proposition}[Tangent bundles and products of real analytic manifolds]
  \label{prop:analytic-bundle}
  \lean{LocallyConvex.Manifolds.realAnalyticOn_tangentChart_transition,
    LocallyConvex.Manifolds.tangentChart_transition_mem_realAnalyticGroupoid,
    LocallyConvex.Manifolds.TangentBundle.realAnalyticManifold,
    LocallyConvex.Manifolds.RealAnalyticOn.prodMap,
    LocallyConvex.Manifolds.prod_mem_realAnalyticGroupoid,
    LocallyConvex.Manifolds.Prod.realAnalyticManifold,
    LocallyConvex.Manifolds.realAnalyticMOn_fst, LocallyConvex.Manifolds.realAnalyticMOn_snd,
    LocallyConvex.Manifolds.RealAnalyticMOn.prodMk}
  \leanok
  \uses{def:analytic-manifold, def:analytic-maps, def:tangent-bundle}
  If $M$ is a real analytic manifold modeled on $E$, its tangent bundle $\T M$ is a real analytic
  manifold modeled on $E \times E$. Products of real analytic manifolds are real analytic
  manifolds, and the projections and pairings of real analytic maps are real analytic.
\end{proposition}

\begin{proof}
  \leanok
  The transition maps of $\T M$ are $(y, w) \mapsto (\tau(y), \dd\tau(y)\, w)$ for transition maps
  $\tau$ of $M$, and the differential of a real analytic map is real analytic: if $g$ is a complex
  analytic extension of $\tau$, then $\dd g$ is one of $\dd\tau$ (lc-analysis). Products of real
  analytic maps are real analytic because products of complex analytic extensions are complex
  analytic extensions.
\end{proof}

\begin{definition}[Real analytic local additions {\cite[p.~325]{GN26}}]
  \label{def:analytic-local-addition}
  \lean{LocallyConvex.Manifolds.LocalAddition.IsRealAnalytic,
    LocallyConvex.Manifolds.LocalAddition.IsRealAnalytic.realAnalyticMOn_toFun}
  \leanok
  \uses{def:local-addition, prop:analytic-bundle}
  A local addition $\Sigma$ on a real analytic manifold $M$ is \emph{real analytic} if
  $\theta = (\pi, \Sigma)$ is a real analytic diffeomorphism onto its open image. Then $\Sigma$ is
  real analytic.
\end{definition}

A real analytic manifold is a smooth manifold, so a local addition on it makes sense, and being
real analytic is a property of a smooth local addition. The real analytic and the smooth maximal
atlases differ; the charted space structure of $\T M$ is built from a real analytic atlas of $M$.

\begin{corollary}[Real analytic local diffeomorphisms {\cite[Cor.~3.8.19, Cor.~3.8.21]{GN26}}]
  \label{cor:local-diffeo-analytic}
  \lean{LocallyConvex.Manifolds.realAnalyticMOn_of_local,
    LocallyConvex.Manifolds.OpenPartialHomeomorph.realAnalyticMOn_of_local,
    LocallyConvex.Manifolds.OpenPartialHomeomorph.realAnalyticMOn_symm_of_local,
    LocallyConvex.Manifolds.exists_realAnalyticMOn_diffeo_of_isCompact,
    LocallyConvex.Manifolds.exists_realAnalyticMOn_diffeo_of_isSigmaCompact}
  \leanok
  \uses{prop:local-diffeo-compact, prop:local-diffeo-sigma-compact, def:analytic-maps}
  Real analyticity of maps between real analytic manifolds is a local property. Hence
  Propositions~\ref{prop:local-diffeo-compact} and~\ref{prop:local-diffeo-sigma-compact} hold for
  real analytic local diffeomorphisms.
\end{corollary}

\begin{proof}
  \leanok
  Locality: real analyticity is checked in charts, and on chart domains it is local in the model
  \cite[Lemma~2.2.3]{GN26}. The rest is the topological part of the two propositions.
\end{proof}

\begin{definition}[Split real analytic submanifolds {\cite[Def.~3.1.10]{GN26}}]
  \label{def:split-analytic-submanifold}
  \lean{LocallyConvex.Manifolds.IsAnalyticSliceChart,
    LocallyConvex.Manifolds.IsSplitAnalyticSubmanifold,
    LocallyConvex.Manifolds.IsSplitAnalyticSubmanifold.toEmbedded,
    LocallyConvex.Manifolds.IsSplitAnalyticSubmanifold.chartedSpace}
  \leanok
  \uses{def:embedded-submanifold, def:analytic-manifold}
  Let $S \subseteq F$ be a complemented subspace, with a continuous linear projection $P \colon F
  \to S$. A subset $X \subseteq F$ is a \emph{split real analytic submanifold} modeled on $S$ if
  every point of $X$ lies in a \emph{real analytic slice chart}: a real analytic diffeomorphism
  $\varphi$ between open subsets of $F$ with $\varphi(U \cap X) = \varphi(U) \cap S$. The
  restrictions of these charts form an atlas of $X$. Such an $X$ is an embedded submanifold.
\end{definition}

\begin{theorem}[Split real analytic submanifolds are real analytic manifolds]
  \label{thm:split-analytic-submanifold}
  \lean{LocallyConvex.Manifolds.IsSplitAnalyticSubmanifold.realAnalyticOn_sliceChart_transition,
    LocallyConvex.Manifolds.IsSplitAnalyticSubmanifold.realAnalyticManifold,
    LocallyConvex.Manifolds.IsSplitAnalyticSubmanifold.realAnalyticMOn_of_val,
    LocallyConvex.Manifolds.IsSplitAnalyticSubmanifold.realAnalyticMOn_val,
    LocallyConvex.Manifolds.IsSplitAnalyticSubmanifold.ambient,
    LocallyConvex.Manifolds.IsSplitAnalyticSubmanifold.realAnalyticMOn_ambient,
    LocallyConvex.Manifolds.IsSplitAnalyticSubmanifold.realAnalyticMOn_of_ambient}
  \leanok
  \uses{def:split-analytic-submanifold, prop:analytic-bundle}
  A split real analytic submanifold $X \subseteq F$ is a real analytic manifold modeled on $S$. A
  map into $X$ is real analytic if and only if it is real analytic into $F$. The ambient map
  $\T X \to F \times F$, $v \mapsto (\pi v, w)$, with $w \in T_p X \subseteq F$ the ambient vector of
  $v$, is real analytic, and a map into $\T X$ is real analytic as soon as its composition with the
  ambient map is.
\end{theorem}

\begin{proof}
  \leanok
  As for smooth embedded submanifolds (Proposition~\ref{prop:ambient-local-addition}), with
  real analytic maps in place of smooth ones. The only new point is corestriction: a map with
  values in $S$ which is real analytic into $F$ is real analytic into $S$ because it equals its
  composition with $P$.
\end{proof}

\begin{remark}[Why complemented]
  For smooth maps, a map with values in a closed subspace $S$ is smooth into $S$. For real analytic
  maps in Milnor's sense this is less clear: a complex analytic extension to $D_{\mathbb{C}}$ of a
  map with values in $S$ need not take values in $S_{\mathbb{C}}$. Assuming $S$ complemented avoids
  the issue. In finite dimensions, and for closed subspaces of Hilbert spaces, this is no
  restriction; for general locally convex (even Banach) spaces it is. The real analytic atlas
  differs from the atlas of all smooth slice charts, so the two tangent bundles of $X$ are
  different types in Lean, although the tangent spaces $T_pX \subseteq F$ agree.
\end{remark}

\begin{proposition}[Real analytic local additions from ambient formulas]
  \label{prop:analytic-ambient-local-addition}
  \lean{LocallyConvex.Manifolds.AnalyticAmbientAddition,
    LocallyConvex.Manifolds.AnalyticAmbientAddition.realAnalyticMOn_theta,
    LocallyConvex.Manifolds.AnalyticAmbientAddition.realAnalyticMOn_thetaInv,
    LocallyConvex.Manifolds.AnalyticAmbientAddition.toLocalAddition,
    LocallyConvex.Manifolds.AnalyticAmbientAddition.toLocalAddition_toFun,
    LocallyConvex.Manifolds.AnalyticAmbientAddition.isRealAnalytic_toLocalAddition}
  \leanok
  \uses{thm:split-analytic-submanifold, def:analytic-local-addition,
    prop:ambient-local-addition}
  Let $X \subseteq F$ be a split real analytic submanifold, and $\sigma$, $\tau$ real analytic maps
  on open subsets of $F \times F$ satisfying the hypotheses of
  Proposition~\ref{prop:ambient-local-addition}. Then $\Sigma(v) = \sigma(\pi v, v)$ is a real
  analytic local addition on $X$.
\end{proposition}

\begin{proof}
  \leanok
  By Theorem~\ref{thm:split-analytic-submanifold}, $\theta = (\pi, \Sigma)$ and
  $\theta^{-1}(p, q) = (p, \tau(p, q))$ are real analytic, because their compositions with the
  inclusion and with the ambient map are. Glöckner and Neeb do not discuss ambient formulas; we do
  not claim this argument is new.
\end{proof}

\section{Manifolds of continuous maps}

Let $K$ be a compact Hausdorff space and $N$ a manifold with a local addition
$\Sigma \colon \Omega \to N$, $\theta = (\pi, \Sigma)$.

\begin{definition}[The spaces $\Gamma_f$ {\cite[4.1.30, Lemma~4.1.31]{GN26}}]
  \label{def:gamma-f}
  \lean{LocallyConvex.Manifolds.coordAt, LocallyConvex.Manifolds.continuousAt_iff_coordAt,
    LocallyConvex.Manifolds.sectionSubmodule, LocallyConvex.Manifolds.Sec,
    LocallyConvex.Manifolds.Sec.CoordIdx, LocallyConvex.Manifolds.Sec.coordFun,
    LocallyConvex.Manifolds.Sec.toPi, LocallyConvex.Manifolds.Sec.instLocallyConvexSpaceReal,
    LocallyConvex.Manifolds.Sec.instT2Space, LocallyConvex.Manifolds.Sec.isClosed_range_toPi,
    LocallyConvex.Manifolds.Sec.isInducing_toTangent}
  \leanok
  \uses{def:tangent-bundle}
  For $f \in C(K, N)$, $\Gamma_f = \{\tau \in C(K, \T N) : \pi \circ \tau = f\}$, with the
  pointwise linear structure and the initial topology with respect to the fiber coordinates
  $\tau \mapsto (x \mapsto \tau(x) \text{ in the chart } i) \in C(L, E)$, for charts $i$ of $N$ and
  compact $L \subseteq K$ with $f(L)$ in the domain of $i$. It is a Hausdorff locally convex space,
  its topology is the one induced by $C(K, \T N)$, and the fiber coordinates embed it as a closed
  subspace of a product.
\end{definition}

\begin{proposition}[Charts of $C(K, N)$ {\cite[proof of Thm.~4.2.2]{GN26}}]
  \label{prop:mapping-charts}
  \lean{LocallyConvex.Manifolds.LocalAddition.chartSource,
    LocallyConvex.Manifolds.LocalAddition.chartTarget, LocallyConvex.Manifolds.LocalAddition.chart,
    LocallyConvex.Manifolds.LocalAddition.mem_chart_source,
    LocallyConvex.Manifolds.LocalAddition.toTangent_chart,
    LocallyConvex.Manifolds.LocalAddition.chart_symm_apply}
  \leanok
  \uses{def:gamma-f, def:local-addition}
  The map $\varphi_f \colon O'_f \to O_f$, $g \mapsto \theta^{-1} \circ (f, g)$, from
  $O'_f = \{g : (f, g)(K) \subseteq \theta(\Omega)\}$ (open in the compact-open topology) onto
  $O_f = \Gamma_f \cap C(K, \Omega)$ (open in $\Gamma_f$) is a homeomorphism, with inverse
  $\tau \mapsto \Sigma \circ \tau$.
\end{proposition}

\begin{proof}
  \leanok
  Postcomposition with maps that are continuous on open sets is continuous in the compact-open
  topology ($K$ is Hausdorff), and $\Gamma_f$ carries the topology induced by $C(K, \T N)$.
\end{proof}

\begin{theorem}[$C(K, N)$ is a manifold {\cite[Thm.~4.2.2(a)--(c), $\ell = 0$]{GN26}}]
  \label{thm:mapping-manifold}
  \lean{LocallyConvex.Analysis.contDiffMBOn_pi,
    LocallyConvex.Analysis.contDiffMBOn_of_isClosedEmbedding,
    LocallyConvex.Analysis.restrictPieces, LocallyConvex.Analysis.isEmbedding_restrictPieces,
    LocallyConvex.Analysis.isClosed_range_restrictPieces,
    LocallyConvex.Manifolds.LocalAddition.Sec.contDiffMBOn_of_coordFun,
    LocallyConvex.Manifolds.LocalAddition.transMap,
    LocallyConvex.Manifolds.LocalAddition.smoothMBOn_transMap,
    LocallyConvex.Manifolds.LocalAddition.transition,
    LocallyConvex.Manifolds.LocalAddition.contDiffMBOn_coordFun_transition_of_adapted,
    LocallyConvex.Manifolds.LocalAddition.smoothMBOn_transition,
    LocallyConvex.Manifolds.VaryingAtlas, LocallyConvex.Manifolds.LocalAddition.mappingAtlas,
    LocallyConvex.Manifolds.LocalAddition.smoothMBMOn_eval_chart_symm}
  \leanok
  \uses{prop:mapping-charts, prop:prod-local-addition}
  The charts $\varphi_f$, $f \in C(K, N)$, form an atlas of $C(K, N)$ with varying models
  $\Gamma_f$ and smooth transition maps; the underlying topology is the compact-open topology,
  and the point evaluations $g \mapsto g(x)$ are smooth.
\end{theorem}

\begin{proof}
  \leanok
  A map into $\Gamma_g$ is smooth if its fiber coordinates are, because these form a closed linear
  embedding into a product. The coordinates over a compact set are determined by those over the
  pieces of a finite cover on which $f$ and $g$ lie in charts $i$ and $j$ of $N$, again through a
  closed embedding. Over such a piece, the transition map $\varphi_g \circ \varphi_f^{-1}$ is the
  superposition operator $\gamma \mapsto h \circ \gamma$ of the smooth map
  $h(y, z, v) = \theta^{-1}(j^{-1} z, \Sigma(i^{-1} y, v))$ (in the chart $j$), applied to a curve
  that depends affinely on the section; superposition operators are smooth by
  \cite[Prop.~1.7.1]{GN26} (formalized in lc-analysis). Glöckner and Neeb deduce the smoothness
  from the exponential law for sections instead.
\end{proof}

\section{Manifolds of \texorpdfstring{$C^\ell$}{C\^{}l} maps}

Let $M$ be a compact smooth manifold, $N$ a manifold with a local addition
$\Sigma \colon \Omega \to N$, $\theta = (\pi, \Sigma)$, and $\ell \in \N_0 \cup \{\infty\}$.

\begin{lemma}[Compact manifolds are finite-dimensional]
  \label{lem:compact-finite-dim}
  \lean{LocallyConvex.Manifolds.locallyCompactSpace_of_compactSpace,
    LocallyConvex.Manifolds.finiteDimensional_of_compactSpace}
  \leanok
  If $M$ is nonempty and compact, its model is locally compact, hence finite-dimensional by
  Riesz's theorem; so $M$ is modeled on some $\R^n$, and only $N$ may be infinite-dimensional.
  The results below use only the local compactness of the model, inside the proofs (the case
  $M = \emptyset$, where $C^\ell(M, N)$ is a point, is treated separately), so that their
  statements carry only the hypothesis that $M$ is compact.
\end{lemma}

\begin{proof}
  \leanok
  A chart sends a compact neighborhood of a point onto a compact neighborhood in the model, and a
  topological group with a compact neighborhood is locally compact.
\end{proof}

\begin{proposition}[Superposition operators on $C^r(U, E)$]
  \label{prop:compLeftOn-smooth}
  \lean{LocallyConvex.Analysis.restrictCompacts,
    LocallyConvex.Analysis.isEmbedding_restrictCompacts,
    LocallyConvex.Analysis.isClosed_range_restrictCompacts,
    LocallyConvex.Analysis.ContDiffMBMap.compLeftOn,
    LocallyConvex.Analysis.ContDiffMBMap.dMap_compLeftOn,
    LocallyConvex.Analysis.ContDiffMBMap.contDiffMBOn_iteratedD_zero_compLeftOn,
    LocallyConvex.Analysis.ContDiffMBMap.contDiffMBOn_compLeftOn}
  \leanok
  Let $Z$ be locally compact, $U \subseteq Z$ open, $r \in \N_0$, and $h \colon E \to F$ smooth on
  an open set $O$. If $g \colon X \to C^r(U, E)$ is $C^k$ on an open set $W$ and
  $g(w)(U) \subseteq O$ for $w \in W$, then $w \mapsto h \circ g(w) \in C^r(U, F)$ is $C^k$ on $W$.
\end{proposition}

\begin{proof}
  \leanok
  Induction on $r$. Since $C^r(U, F)$ is a closed subspace of $\prod_{j \le r} C(U \times Z^j, F)$
  \cite[Lemma~1.7.15]{GN26}, it suffices to treat each $d^{(j)}$. For $j = 0$, restriction to the
  compact subsets of $U$ (a closed embedding, $Z$ being locally compact) reduces the claim to
  superposition operators on $C(C, E)$, $C$ compact, which are smooth
  \cite[Prop.~1.7.1]{GN26}. For $j \ge 1$, $d^{(j)}(h \circ \gamma)$ depends continuously and
  linearly on $d(h \circ \gamma) = dh \circ (\gamma \circ \mathrm{pr}_1, d\gamma)$, a superposition
  by $dh$ on $C^{r-1}(U \times Z, E \times E)$, and the induction hypothesis applies. Glöckner
  and Neeb obtain the smoothness they need from the exponential law
  \cite[Prop.~4.1.34]{GN26} instead.
\end{proof}

\begin{lemma}[Restriction to open covers]
  \label{lem:restrict-cover}
  \lean{LocallyConvex.Analysis.ContDiffMBMap.restrictCLM,
    LocallyConvex.Analysis.ContDiffMBMap.restrictCover,
    LocallyConvex.Analysis.ContDiffMBMap.tendsto_of_tendsto_restrictCLM,
    LocallyConvex.Analysis.ContDiffMBMap.isEmbedding_restrictCover,
    LocallyConvex.Analysis.ContDiffMBMap.isClosed_range_restrictCover,
    LocallyConvex.Analysis.ContDiffMBMap.contDiffMBOn_of_restrictCLM}
  \leanok
  For an open cover $(V_a)$ of $U$, the restrictions $C^r(U, F) \to \prod_a C^r(V_a, F)$ form a
  linear topological embedding with closed range (compare \cite[Lemma~1.7.43, Lemma~4.1.9]{GN26}).
\end{lemma}

\begin{definition}[Jets on compact sets in charts; the topology of $C^\ell(M, N)$]
  \label{def:jet-topology}
  \lean{LocallyConvex.Manifolds.jetOn, LocallyConvex.Manifolds.jetOn_apply,
    LocallyConvex.Manifolds.jetOn_mono, LocallyConvex.Manifolds.CkMap, LocallyConvex.Manifolds.CkMap.CompactIdx,
    LocallyConvex.Manifolds.CkMap.CompactIdx.jet, LocallyConvex.Manifolds.CkMap.jetTopology,
    LocallyConvex.Manifolds.CkMap.instTopologicalSpace,
    LocallyConvex.Manifolds.CkMap.continuous_jetTopology_iff,
    LocallyConvex.Manifolds.SmoothMap, LocallyConvex.Manifolds.SmoothMap.jetTopologySm,
    LocallyConvex.Manifolds.SmoothMap.instTopologicalSpace}
  \leanok
  For charts $j$ of $M$ and $i$ of $N$, a compact set $K$ in the target of $j$ (no interior is
  required) and $m \le \ell$, the \emph{$m$-jet on $K$} of $g \in C^\ell(M, N)$ is
  $\dd^{(m)}(i \circ g \circ j^{-1})|_{K \times H^m} \in C(K \times H^m, E)$, defined on the open set
  of maps with $g(j^{-1}(K))$ in the domain of $i$. The \emph{compact-open $C^\ell$ topology} of
  $C^\ell(M, N)$, or \emph{jet topology}, is the coarsest topology finer than the compact-open
  topology of $C(M, N)$ that makes all the jets continuous, with the compact-open topology on
  $C(K \times H^m, E)$. A map into $C^\ell(M, N)$ is continuous if and only if it is continuous
  into $C(M, N)$ and its jets are continuous where defined. For $\ell = \infty$, $C^\infty(M, N)$
  carries the initial topology with respect to the inclusions into the $C^\ell(M, N)$
  \cite[Remark~4.1.3(b)]{GN26}. This definition makes sense for every manifold $M$ (compact or
  not) modeled on any Hausdorff locally convex space $H$; for Hausdorff $M$ it is the topology of
  \cite[Def.~4.1.2]{GN26} (Theorem~\ref{thm:jet-topology-gn}).
\end{definition}

This is the weak (compact-open) $C^r$ topology of \cite[Ch.~2, \S1, pp.~34--35]{Hir76}, with the
uniform bounds on the derivatives over $K$ replaced by the compact-open topology of
$C(K \times H^m, E)$, as in the compact-open $C^k$ topology of \cite[Def.~1.7.10]{GN26} for maps on
open sets.

\begin{definition}[The chart topology]
  \label{def:chart-topology}
  \lean{LocallyConvex.Manifolds.CkMap.ChartIdx, LocallyConvex.Manifolds.CkMap.ChartIdx.rep,
    LocallyConvex.Manifolds.CkMap.chartTopology, LocallyConvex.Manifolds.SmoothMap.chartTopologySm,
    LocallyConvex.Manifolds.CkMap.continuous_of_rep,
    LocallyConvex.Manifolds.CkMap.continuousOn_of_rep}
  \leanok
  \uses{def:jet-topology}
  The \emph{chart topology} of $C^\ell(M, N)$ is the coarsest topology finer than the
  compact-open topology for which the local representatives
  $g \mapsto i \circ g \circ j^{-1} \in C^\ell(V, E)$ are continuous on the open sets
  $\{g : g(j^{-1}(L)) \subseteq \mathrm{dom}\, i\}$, for charts $j$ of $M$, $i$ of $N$, compact
  $L$ and open $V \subseteq L$; for $\ell = \infty$, the initial topology with respect to the finite
  orders. It is not the topology of $C^\ell(M, N)$: if the model of $M$ is not locally compact,
  compact sets have empty interior, every such $V$ is empty, and the chart topology collapses to
  the compact-open $C^0$ topology.
\end{definition}

\begin{proposition}[The chart topology and the jet topology]
  \label{prop:chart-jet}
  \lean{LocallyConvex.Manifolds.CkMap.ChartIdx.jetOn_rep,
    LocallyConvex.Manifolds.CkMap.jetTopology_le_chartTopology,
    LocallyConvex.Manifolds.CkMap.chartTopology_le_jetTopology,
    LocallyConvex.Manifolds.CkMap.chartTopology_eq_jetTopology,
    LocallyConvex.Manifolds.CkMap.topology_eq_chartTopology,
    LocallyConvex.Manifolds.CkMap.topology_eq_chartTopology_of_compactSpace,
    LocallyConvex.Manifolds.SmoothMap.chartTopologySm_eq_jetTopologySm,
    LocallyConvex.Manifolds.SmoothMap.topology_eq_chartTopologySm_of_compactSpace}
  \leanok
  \uses{def:jet-topology, def:chart-topology, lem:compact-finite-dim}
  The jet topology is always finer than the chart topology. If the model of $M$ is locally
  compact, in particular if $M$ is compact, the two topologies coincide,
  $\ell \in \N_0 \cup \{\infty\}$.
\end{proposition}

\begin{proof}
  \leanok
  On a compact set $K \subseteq V$ the jets of the local representative on $V$ are the jets of $g$
  on $K$, and the compact-open topology of $C(V \times H^m, E)$ is determined by the compact
  subsets of $V \times H^m$; this gives the first inclusion. Conversely, if the model is locally
  compact, a compact set $K$ in the open domain of $i \circ g_0 \circ j^{-1}$ lies in the interior
  $V$ of a compact set $L$ in that domain; for $g$ near $g_0$ the jets of $g$ on $K$ are
  restrictions of the differentials of the local representative on $V$. A nonempty compact
  manifold has a locally compact model (Lemma~\ref{lem:compact-finite-dim}).
\end{proof}

\begin{definition}[$\Gamma_f$ {\cite[4.1.30, Lemma~4.1.31]{GN26}}]
  \label{def:gamma-f-ck}
  \lean{LocallyConvex.Manifolds.chartDomOpens,
    LocallyConvex.Manifolds.localCoord, LocallyConvex.Manifolds.sectionSubmoduleCk,
    LocallyConvex.Manifolds.SecCk, LocallyConvex.Manifolds.SecCk.coordFun,
    LocallyConvex.Manifolds.SecCk.toPi, LocallyConvex.Manifolds.SecCk.instLocallyConvexSpaceReal,
    LocallyConvex.Manifolds.SecCk.instT2Space, LocallyConvex.Manifolds.SecCk.isClosed_range_toPi,
    LocallyConvex.Manifolds.SecCk.contDiffMBOn_of_coordFun, LocallyConvex.Manifolds.SecCk.toSec,
    LocallyConvex.Manifolds.SecSm, LocallyConvex.Manifolds.SecSm.isClosed_range_toPi}
  \leanok
  \uses{def:gamma-f, def:jet-topology}
  The space $\Gamma_f$ of $C^\ell$ sections of $TN$ along $f$ carries the initial topology with
  respect to the local representatives in charts; it is a Hausdorff locally convex space,
  embedded as a closed subspace of a product of spaces $C^\ell(D, E)$. For $\ell = \infty$ it
  carries the initial topology with respect to the finite orders \cite[Remark~4.1.3(b)]{GN26}.
  Glöckner and Neeb define this topology through $C^\ell(M, TN)$; the chart description used here
  corresponds to their Definition~4.1.16, Proposition~4.1.26 and the proof of
  Proposition~4.1.34.
\end{definition}

\begin{theorem}[$C^\ell(M, N)$ is a manifold {\cite[Thm.~4.2.2(a)--(c)]{GN26}}]
  \label{thm:mapping-manifold-ck}
  \lean{LocallyConvex.Manifolds.LocalAddition.chartLocal,
    LocallyConvex.Manifolds.LocalAddition.chartLocal_spec,
    LocallyConvex.Manifolds.LocalAddition.invLocal,
    LocallyConvex.Manifolds.LocalAddition.invLocal_spec,
    LocallyConvex.Manifolds.LocalAddition.transMap_spec,
    LocallyConvex.Manifolds.LocalAddition.continuousOn_chartFunCk,
    LocallyConvex.Manifolds.LocalAddition.continuousOn_chartInvFunCk,
    LocallyConvex.Manifolds.LocalAddition.chartCk,
    LocallyConvex.Manifolds.LocalAddition.smoothMBOn_transitionCk,
    LocallyConvex.Manifolds.LocalAddition.mappingAtlasCk,
    LocallyConvex.Manifolds.LocalAddition.smoothMBMOn_eval_chartCk_symm,
    LocallyConvex.Manifolds.LocalAddition.chartSm,
    LocallyConvex.Manifolds.LocalAddition.smoothMBOn_transitionSm,
    LocallyConvex.Manifolds.LocalAddition.mappingAtlasSm,
    LocallyConvex.Manifolds.LocalAddition.smoothMBMOn_eval_chartSm_symm}
  \leanok
  \uses{def:gamma-f-ck, prop:chart-jet, prop:compLeftOn-smooth, lem:restrict-cover,
    thm:mapping-manifold}
  The charts $\varphi_f \colon g \mapsto \theta^{-1} \circ (f, g)$, $f \in C^\ell(M, N)$, are
  partial homeomorphisms for the compact-open $C^\ell$ topology onto open subsets of $\Gamma_f$,
  with inverses $\tau \mapsto \Sigma \circ \tau$; they form an atlas of $C^\ell(M, N)$ with varying
  models $\Gamma_f$ and smooth transition maps, and the point evaluations are smooth.
\end{theorem}

\begin{proof}
  \leanok
  Near each point, $f$ and $g$ lie in charts $i$ and $j$ of $N$, and the local representative of
  $\varphi_g \circ \varphi_f^{-1}$ is the superposition operator of the smooth map
  $(y, z, v) \mapsto \theta^{-1}(j^{-1} z, \Sigma(i^{-1} y, v))$ applied to
  $(i \circ f, j \circ g, \tau)$, which depends affinely on $\tau$
  (Proposition~\ref{prop:compLeftOn-smooth}); Lemma~\ref{lem:restrict-cover} glues the local
  statements. The continuity of the charts is proved the same way, over compact neighborhoods in
  charts of $M$, in the chart topology, which is the topology of $C^\ell(M, N)$ for compact $M$
  (Proposition~\ref{prop:chart-jet}). The case $\ell = \infty$ follows from the finite orders.
\end{proof}

\section{Smooth maps between atlases with varying models}

\begin{definition}[$C^k$ maps for atlases with varying models {\cite[Def.~3.2.13]{GN26}}]
  \label{def:varying-ck}
  \lean{LocallyConvex.Manifolds.VaryingAtlas.ContMDiffOn,
    LocallyConvex.Manifolds.VaryingAtlas.SmoothOn, LocallyConvex.Manifolds.VaryingAtlas.Smooth,
    LocallyConvex.Manifolds.VaryingAtlas.ofChartedSpace}
  \leanok
  \uses{thm:mapping-manifold}
  A map $f \colon X \to Y$ between spaces with atlases with varying models is $C^k$ on an open set
  $U$ if it is continuous on $U$ and all its local representatives $\psi_j \circ f \circ
  \varphi_i^{-1}$ are $C^k$ on their (open) domains. A manifold modeled on a single space $G$ is an
  atlas with varying models through all the charts of its atlas.
\end{definition}

\begin{proposition}[Basic properties]
  \label{prop:varying-ck-basic}
  \lean{LocallyConvex.Manifolds.VaryingAtlas.contMDiffOn_of_local,
    LocallyConvex.Manifolds.VaryingAtlas.contMDiffOn_iff_local,
    LocallyConvex.Manifolds.VaryingAtlas.ContMDiffOn.mono,
    LocallyConvex.Manifolds.VaryingAtlas.ContMDiffOn.congr,
    LocallyConvex.Manifolds.VaryingAtlas.ContMDiffOn.comp,
    LocallyConvex.Manifolds.VaryingAtlas.smoothOn_id,
    LocallyConvex.Manifolds.VaryingAtlas.contMDiffOn_ofChartedSpace_iff,
    LocallyConvex.Manifolds.VaryingAtlas.smoothOn_chart,
    LocallyConvex.Manifolds.VaryingAtlas.smoothOn_chart_symm}
  \leanok
  \uses{def:varying-ck}
  Being $C^k$ can be checked in one pair of charts around each point; $C^k$ maps compose, the
  identity is smooth, and restrictions of $C^k$ maps to open subsets are $C^k$. For manifolds with a
  single model these are the $C^k$ maps of Definition~\ref{def:contmdiff}. Charts are smooth,
  with smooth inverses.
\end{proposition}

\begin{proof}
  \leanok
  Near each point, a local representative in arbitrary charts is the composition of the given one
  with transition maps, which are smooth.
\end{proof}

\begin{proposition}[Equivalent atlases]
  \label{prop:varying-equivalent}
  \lean{LocallyConvex.Manifolds.VaryingAtlas.Equivalent,
    LocallyConvex.Manifolds.VaryingAtlas.Compatible,
    LocallyConvex.Manifolds.VaryingAtlas.smooth_id_iff,
    LocallyConvex.Manifolds.VaryingAtlas.equivalent_iff_compatible,
    LocallyConvex.Manifolds.VaryingAtlas.union}
  \leanok
  \uses{prop:varying-ck-basic}
  Two atlases with varying models on the same space define the same smooth structure (the identity
  is smooth in both directions) if and only if all the transition maps between their charts are
  smooth, that is, if and only if their union is an atlas.
\end{proposition}

\begin{proof}
  \leanok
  The local representatives of the identity are the transition maps.
\end{proof}

\begin{definition}[Tangent spaces and tangent maps]
  \label{def:varying-tangent}
  \lean{LocallyConvex.Manifolds.VaryingAtlas.TangentSpace,
    LocallyConvex.Manifolds.VaryingAtlas.tangentCoordChange,
    LocallyConvex.Manifolds.VaryingAtlas.tangentCoordChange_self,
    LocallyConvex.Manifolds.VaryingAtlas.tangentCoordChange_comp,
    LocallyConvex.Manifolds.VaryingAtlas.tangentCoordChangeEquiv,
    LocallyConvex.Manifolds.VaryingAtlas.tangentMapAt,
    LocallyConvex.Manifolds.VaryingAtlas.dMB_chart_eq_tangentMapAt}
  \leanok
  \uses{prop:varying-ck-basic}
  The tangent space at $x$ is the model of the preferred chart at $x$; the derivatives of the
  transition maps at the point are isomorphisms of topological vector spaces satisfying the
  cocycle condition. The tangent map $T_x f$ of a $C^1$ map is the derivative of its local
  representative in the preferred charts; in other charts it is conjugated by the changes of
  tangent coordinates. (In infinite dimensions this kinematic tangent space need not agree with
  the space of derivations.)
\end{definition}

\section{\texorpdfstring{$C^{k,\ell}$}{C\^{}\{k,l\}} maps on products of manifolds}

\begin{proposition}[Local properties and the chain rule]
  \label{prop:ckl-chain}
  \lean{LocallyConvex.Analysis.ContDiffMBOn₂.mono, LocallyConvex.Analysis.ContDiffMBOn₂.of_local,
    LocallyConvex.Analysis.ContDiffMBOn₂.exists_curry,
    LocallyConvex.Analysis.contDiffMBOn₂_of_curry,
    LocallyConvex.Analysis.ContDiffMBOn₂.comp_prodMap}
  \leanok
  \uses{prop:compLeftOn-smooth}
  Being $C^{k,\ell}$ on products of open sets is a local property. If $f$ is $C^{k,\ell}$ on
  $U \times V$, $\alpha$ is $C^k$, $\beta$ is $C^\ell$ and $\gamma$ is smooth (with the obvious
  conditions on the images), and the second factor of the domain of $\beta$ is locally compact,
  then $\gamma \circ f \circ (\alpha \times \beta)$ is $C^{k,\ell}$
  (compare \cite[Props.~1.7.26, 1.7.28]{GN26}).
\end{proposition}

\begin{proof}
  \leanok
  By the exponential law \cite[Thm.~1.7.30]{GN26}, $x \mapsto f(x, \cdot)$ is $C^k$ into
  $C^\ell(V, F)$; pulling back by $\beta$ is continuous linear \cite[Prop.~1.7.11]{GN26},
  superposition by $\gamma$ is smooth (Proposition~\ref{prop:compLeftOn-smooth}), and uncurrying
  needs local compactness. Glöckner and Neeb use Faà di Bruno's formula instead.
\end{proof}

\begin{definition}[$C^{k,\ell}$ maps {\cite[Def.~3.5.24]{GN26}}]
  \label{def:ckl-manifold}
  \lean{LocallyConvex.Manifolds.localRep₂, LocallyConvex.Manifolds.IsChartBox,
    LocallyConvex.Manifolds.ContMDiffMB₂, LocallyConvex.Manifolds.ContMDiffMB₂.continuous,
    LocallyConvex.Manifolds.ContMDiffMB₂.contDiffMBOn₂}
  \leanok
  \uses{prop:ckl-chain, def:contmdiff, def:varying-ck}
  Let $L$ carry an atlas with varying models (for instance a manifold)
  \cite[Remark~3.5.28]{GN26}, and $M$, $N$ be manifolds without boundary (manifolds with rough
  boundary are out of scope). A map $f \colon L \times M \to N$ is $C^{k,\ell}$ if around
  each point there are charts $e$, $e'$, $c$ and open sets $U \ni e(x)$, $V \ni e'(y)$ with
  $f(e^{-1}U \times e'^{-1} V) \subseteq \mathrm{dom}\, c$ and
  $c \circ f \circ (e^{-1} \times e'^{-1})$ of class $C^{k,\ell}$ on $U \times V$. Such maps are
  continuous, and, if the model of $M$ is locally compact, their local representatives on all such
  boxes, in all charts, are $C^{k,\ell}$ \cite[Remark~3.5.26(a)]{GN26} (this holds when $M$ is
  compact, Lemma~\ref{lem:compact-finite-dim}).
\end{definition}

\section{The exponential law for manifolds of mappings}

\begin{theorem}[Exponential law {\cite[Thm.~4.2.2(a)]{GN26}}]
  \label{thm:mapping-exp-law}
  \lean{LocallyConvex.Manifolds.LocalAddition.contMDiffMB₂_of_contMDiffOn,
    LocallyConvex.Manifolds.LocalAddition.continuous_of_contMDiffMB₂,
    LocallyConvex.Manifolds.LocalAddition.contMDiffOn_of_contMDiffMB₂,
    LocallyConvex.Manifolds.LocalAddition.contMDiffOn_mappingAtlasCk_iff,
    LocallyConvex.Manifolds.LocalAddition.smooth_toCk,
    LocallyConvex.Manifolds.LocalAddition.contMDiffOn_mappingAtlasSm_iff}
  \leanok
  \uses{thm:mapping-manifold-ck, def:ckl-manifold, prop:varying-ck-basic}
  Let $M$ be compact (without boundary) and $N$ admit a local addition. For every space $L$ with
  an atlas with varying models and every $k \in \N_0$, a map $g \colon L \to C^\ell(M, N)$ is
  $C^k$ if and only if $g^\wedge \colon L \times M \to N$, $(x, y) \mapsto g(x)(y)$, is
  $C^{k,\ell}$ (for $\ell = \infty$: $C^{k,\ell'}$ for every $\ell' \in \N_0$).
\end{theorem}

\begin{proof}
  \leanok
  Locally, $g^\wedge$ reads $(a, b) \mapsto c(\Sigma(\tau(a)(b)))$ with $\tau = \varphi_f \circ g$;
  as a function of $a$ this is a superposition operator applied to a $C^k$ curve in a space of
  $C^\ell$ maps, and uncurrying gives $C^{k,\ell}$. Conversely, the local representatives of
  $\varphi_f \circ g$ are superposition operators applied to the curried local representatives
  of $g^\wedge$, over a finite cover of $M$ by compact sets on which $g(x)$ stays in single charts
  of $N$. Glöckner and Neeb argue through the exponential law for sections
  \cite[Prop.~4.1.34]{GN26}. (Part (b) of \cite[Thm.~4.2.2]{GN26} holds by construction for the
  chart description of the compact-open $C^\ell$ topology; its agreement with
  \cite[Def.~4.1.2]{GN26} is Theorem~\ref{thm:mapping-topology-gn}.)
\end{proof}

\begin{corollary}[The evaluation map {\cite[Thm.~4.2.2(c)]{GN26}}]
  \label{cor:mapping-eval}
  \lean{LocallyConvex.Manifolds.ContMDiffMB₂.contMDiffOn_left,
    LocallyConvex.Manifolds.LocalAddition.contMDiffMB₂_eval,
    LocallyConvex.Manifolds.LocalAddition.smooth_eval,
    LocallyConvex.Manifolds.LocalAddition.contMDiffMB₂_evalSm,
    LocallyConvex.Manifolds.LocalAddition.smooth_evalSm}
  \leanok
  \uses{thm:mapping-exp-law}
  The evaluation map $\varepsilon \colon C^\ell(M, N) \times M \to N$ is $C^{k,\ell}$ for every
  $k$, and the point evaluations $\varepsilon_x$ are smooth.
\end{corollary}

\begin{proof}
  \leanok
  Apply the exponential law to the identity; partial maps of $C^{k,\ell}$ maps are $C^k$.
\end{proof}

\begin{corollary}[Uniqueness {\cite[Thm.~4.2.2(a)]{GN26}}]
  \label{cor:mapping-uniqueness}
  \lean{LocallyConvex.Manifolds.LocalAddition.equivalent_mappingAtlasCk_of_expLaw,
    LocallyConvex.Manifolds.LocalAddition.equivalent_mappingAtlasSm_of_expLaw}
  \leanok
  \uses{thm:mapping-exp-law, prop:varying-equivalent}
  An atlas with varying models on $C^\ell(M, N)$ satisfying the exponential law defines the same
  smooth structure as the atlas of the charts $\varphi_f$.
\end{corollary}

\begin{proof}
  \leanok
  Apply the exponential law to the identity in both directions: the evaluation map is
  $C^{k,\ell}$ for both atlases.
\end{proof}

\begin{theorem}[Tangent spaces of $C^\ell(M, N)$ {\cite[Thm.~4.2.2(d)]{GN26}}]
  \label{thm:mapping-tangent}
  \lean{LocallyConvex.Manifolds.LocalAddition.chartCk_self,
    LocallyConvex.Manifolds.LocalAddition.exists_smoothMBOn_chart_toFun,
    LocallyConvex.Manifolds.LocalAddition.tangentMapAt_eval,
    LocallyConvex.Manifolds.LocalAddition.tangentEquivSecCk,
    LocallyConvex.Manifolds.LocalAddition.coe_tangentEquivSecCk,
    LocallyConvex.Manifolds.LocalAddition.tangentMapAt_evalSm,
    LocallyConvex.Manifolds.LocalAddition.tangentEquivSecSm,
    LocallyConvex.Manifolds.LocalAddition.coe_tangentEquivSecSm,
    LocallyConvex.Manifolds.LocalAddition.mappingAtlasCk_equivalent,
    LocallyConvex.Manifolds.LocalAddition.mappingAtlasSm_equivalent,
    LocallyConvex.Manifolds.LocalAddition.tangentEquivSecCkOf,
    LocallyConvex.Manifolds.LocalAddition.coe_tangentEquivSecCkOf,
    LocallyConvex.Manifolds.LocalAddition.tangentEquivSecSmOf,
    LocallyConvex.Manifolds.LocalAddition.coe_tangentEquivSecSmOf}
  \leanok
  \uses{def:varying-tangent, cor:mapping-eval, cor:mapping-uniqueness, def:local-addition,
    prop:normalize-local-addition}
  For a local addition and $f \in C^\ell(M, N)$, the map
  $\Theta \colon T_f C^\ell(M, N) \to \Gamma_f$, $v \mapsto (T_f \varepsilon_x(v))_{x \in M}$, is an
  isomorphism of topological vector spaces. The smooth structure does not depend on the local
  addition.
\end{theorem}

\begin{proof}
  \leanok
  The chart $\varphi_f$ sends $f$ to the zero section, and in this chart $\varepsilon_x$ reads
  $\tau \mapsto \Sigma(\tau(x))$; by the chain rule and the normalization (3.55),
  $T_f \varepsilon_x(v) = v(x)$, so $\Theta$ is the identity of $\Gamma_f$, the model of the chart
  at $f$, when the local addition is normalized. Glöckner and Neeb compute with geometric tangent
  vectors and assume the local addition normalized, which is no loss of generality
  (Proposition~\ref{prop:normalize-local-addition}). For an arbitrary local addition $\Sigma$, the
  atlases of $\Sigma$ and of its normalization are compatible, and $\Theta$ is the change of
  tangent coordinates between their charts at $f$; it is again $v \mapsto (T_f
  \varepsilon_x(v))_x$, a formula that involves no chart of $C^\ell(M, N)$.
\end{proof}

\section{The compact-open \texorpdfstring{$C^\ell$}{C\^{}l} topology via iterated tangent maps}

\begin{definition}[Iterated tangent bundles; {\cite[Def.~4.1.2]{GN26}}]
  \label{def:gn-topology}
  \lean{LocallyConvex.Manifolds.MBMfld, LocallyConvex.Manifolds.MBMfld.tangent,
    LocallyConvex.Manifolds.MBMfld.iter, LocallyConvex.Manifolds.MBMfld.iterMap,
    LocallyConvex.Manifolds.MBMfld.contMDiffMB_iterMap, LocallyConvex.Manifolds.MBMfld.iterMapC,
    LocallyConvex.Manifolds.MBMfld.gnTopology, LocallyConvex.Manifolds.MBMfld.gnTopologySm,
    LocallyConvex.Manifolds.MBMfld.gnTopology_succ, LocallyConvex.Manifolds.MBMfld.gnTopologySm_eq}
  \leanok
  \uses{def:tangent-map}
  The iterated tangent bundles $T^j M$ and tangent maps $T^j f$ are defined recursively. The
  compact-open $C^\ell$ topology on $C^\ell(M, N)$ is the initial topology with respect to
  $f \mapsto T^j f \in C(T^j M, T^j N)$, $j \le \ell$, compact-open. It is initial with respect to
  $f \mapsto f$ and $f \mapsto Tf \in C^{\ell-1}(TM, TN)$ (as in the proof of
  \cite[Lemma~4.1.6]{GN26}), and for $\ell = \infty$ it is the projective limit of the finite
  orders \cite[Remark~4.1.3(b)]{GN26}.
\end{definition}

\begin{proposition}[Lemma 4.1.6 for the chart topology {\cite[Lemma~4.1.6]{GN26}}]
  \label{prop:chart-topology-succ}
  \lean{LocallyConvex.Manifolds.exists_eq_tangentChart,
    LocallyConvex.Manifolds.continuousOn_iteratedD_zero_rep,
    LocallyConvex.Manifolds.continuous_toC_of_rep, LocallyConvex.Manifolds.tangentIdx_rep_apply,
    LocallyConvex.Manifolds.continuous_of_tangentCk, LocallyConvex.Manifolds.continuous_tangentCk,
    LocallyConvex.Manifolds.chartTopology_zero, LocallyConvex.Manifolds.chartTopology_succ}
  \leanok
  \uses{def:chart-topology, def:gn-topology, prop:compLeftOn-smooth, lem:restrict-cover}
  If the model of $M$ is locally compact, the chart topology of $C^0(M, N)$ is the compact-open
  topology, and that of $C^{\ell+1}(M, N)$ is initial with respect to $f \mapsto f \in C(M, N)$
  and $f \mapsto Tf \in C^\ell(TM, TN)$ (chart topology). (The chart topology is meant for locally
  compact models: otherwise it collapses to the $C^0$ topology.)
\end{proposition}

\begin{proof}
  \leanok
  In self tangent charts, $d(\mathrm{rep}\,f)$ is a component of a local representative of $Tf$;
  in general tangent charts the local representative of $Tf$ is $(F, dF)$ pulled back along a
  smooth map and composed with a smooth map. Localizations use compact neighborhoods.
\end{proof}

\begin{theorem}[The chart topology and Definition 4.1.2 {\cite[Def.~4.1.2]{GN26}}]
  \label{thm:mapping-topology-gn}
  \lean{LocallyConvex.Manifolds.chartTopology_eq_gnTopology,
    LocallyConvex.Manifolds.chartTopologySm_eq_gnTopologySm}
  \leanok
  \uses{prop:chart-topology-succ, def:gn-topology}
  If the model of $M$ is locally compact, the chart topology of $C^\ell(M, N)$ is the compact-open
  $C^\ell$ topology of \cite[Def.~4.1.2]{GN26}, for $\ell \in \N_0 \cup \{\infty\}$.
\end{theorem}

\begin{proof}
  \leanok
  Both topologies satisfy the recursion of \cite[Lemma~4.1.6]{GN26}; induction on $\ell$, and the
  projective limit for $\ell = \infty$.
\end{proof}


\section{The jet topology and Definition 4.1.2}

The topology of $C^\ell(M, N)$ is the jet topology (Definition~\ref{def:jet-topology}), which
works for every Hausdorff locally convex model and every manifold $M$, compact or not. We compare
it with the topology of \cite[Def.~4.1.2]{GN26}.

\begin{lemma}[Localizing compact sets]
  \label{lem:compact-localizable}
  \lean{LocallyConvex.Manifolds.CompactLocalizable,
    LocallyConvex.Manifolds.CompactLocalizable.of_locallyCompactSpace,
    LocallyConvex.Manifolds.CompactLocalizable.of_r1Space,
    LocallyConvex.Manifolds.MBMfld.t2Space_iter,
    LocallyConvex.Manifolds.MBMfld.compactLocalizable_iter_of_t2Space,
    LocallyConvex.Manifolds.MBMfld.compactLocalizable_iter_of_locallyCompactSpace}
  \leanok
  A space is \emph{compact-localizable} if every compact set is covered by finitely many compact
  sets inside prescribed neighborhoods of its points. Hausdorff spaces and locally compact spaces
  are compact-localizable. Tangent bundles of Hausdorff manifolds are Hausdorff, so the iterated
  tangent bundles of a Hausdorff manifold, or of a manifold with a locally compact model, are
  compact-localizable.
\end{lemma}

\begin{proof}
  \leanok
  For Hausdorff spaces this is the finite compact cover lemma used by Glöckner and Neeb
  \cite[Lemma~A.6.11]{GN26}.
\end{proof}

\begin{theorem}[The jet topology is the compact-open $C^\ell$ topology {\cite[Def.~4.1.2]{GN26}}]
  \label{thm:jet-topology-gn}
  \lean{LocallyConvex.Manifolds.jetOn_succ, LocallyConvex.Manifolds.jetOn_dMB,
    LocallyConvex.Manifolds.continuousOn_jetOn_of_isCompact,
    LocallyConvex.Manifolds.tangentJetIdx_jet,
    LocallyConvex.Manifolds.MBMfld.gnTopology_le_jetTopology,
    LocallyConvex.Manifolds.MBMfld.jetTopology_le_gnTopology,
    LocallyConvex.Manifolds.MBMfld.jetTopology_eq_gnTopology,
    LocallyConvex.Manifolds.MBMfld.jetTopology_eq_gnTopology_of_locallyCompactSpace,
    LocallyConvex.Manifolds.MBMfld.jetTopologySm_eq_gnTopologySm,
    LocallyConvex.Manifolds.MBMfld.jetTopologySm_eq_gnTopologySm_of_locallyCompactSpace}
  \leanok
  \uses{def:jet-topology, def:gn-topology, lem:compact-localizable}
  The compact-open $C^\ell$ topology of \cite[Def.~4.1.2]{GN26} is always finer than the jet
  topology. For every Hausdorff manifold $M$, compact or not, modeled on any Hausdorff locally
  convex space, and for every $M$ with a locally compact model, the two topologies coincide,
  $\ell \in \N_0 \cup \{\infty\}$. In particular, for $M$ compact the topology underlying the
  atlases of $C^\ell(M, N)$ and $C^\infty(M, N)$ is that of \cite[Def.~4.1.2]{GN26}
  \cite[Thm.~4.2.2(b)]{GN26}.
\end{theorem}

\begin{proof}
  \leanok
  Induction on the order through the tangent map, as in the proof of \cite[Lemma~4.1.6]{GN26}. In
  the tangent charts the local representative of $\T g$ is $(y, w) \mapsto (F(y), \dd F(y)\,w)$,
  $F = i \circ g \circ j^{-1}$, and the jets of $\dd F$ and those of $F$ determine each other by
  $\dd^{(m+1)}F(x, v_0, v) = \dd^{(m)}(\dd F)((x, v_0), (v, 0))$ \cite[(1.36)]{GN26} and the formula
  for $\dd^{(m)}(\dd F)$, the identities behind \cite[Lemma~4.1.12]{GN26}. Continuity of the jets
  of $\dd F$ on the non-compact set $K \times H$ is reduced to compact subsets. Recovering the
  compact-open topology of $C(\T^n M, \T^n N)$ from compact sets inside charts needs
  Lemma~\ref{lem:compact-localizable}; for non-Hausdorff manifolds with non-locally compact models
  only one inclusion is proved. We do not claim these arguments are new.
\end{proof}

\begin{corollary}[The jet topology and the chart topology, via Definition 4.1.2]
  \label{cor:jet-topology-chart}
  \lean{LocallyConvex.Manifolds.MBMfld.jetTopology_eq_chartTopology}
  \leanok
  \uses{thm:jet-topology-gn, thm:mapping-topology-gn, prop:chart-jet}
  If the model of $M$ is locally compact, the jet topology is the chart topology.
\end{corollary}

\begin{proof}
  \leanok
  This is Proposition~\ref{prop:chart-jet}; for a locally compact model it also follows from
  Theorems~\ref{thm:jet-topology-gn} and~\ref{thm:mapping-topology-gn}, both topologies being that
  of \cite[Def.~4.1.2]{GN26}.
\end{proof}

\begin{theorem}[Targets in any universe]
  \label{thm:mapping-topology-ulift}
  \lean{LocallyConvex.Manifolds.chartedSpaceULift,
    LocallyConvex.Manifolds.mbManifoldULift,
    LocallyConvex.Manifolds.CkMap.uliftHomeomorph,
    LocallyConvex.Manifolds.SmoothMap.uliftHomeomorph,
    LocallyConvex.Manifolds.MBMfld.iterMap_comp,
    LocallyConvex.Manifolds.MBMfld.gnTopology_eq_induced,
    LocallyConvex.Manifolds.MBMfld.gnTopologySm_eq_induced,
    LocallyConvex.Manifolds.CkMap.gnTopology,
    LocallyConvex.Manifolds.SmoothMap.gnTopologySm,
    LocallyConvex.Manifolds.CkMap.gnTopology_eq,
    LocallyConvex.Manifolds.SmoothMap.gnTopologySm_eq,
    LocallyConvex.Manifolds.CkMap.jetTopology_eq_gnTopology,
    LocallyConvex.Manifolds.CkMap.jetTopology_eq_gnTopology_of_locallyCompactSpace,
    LocallyConvex.Manifolds.SmoothMap.jetTopologySm_eq_gnTopologySm,
    LocallyConvex.Manifolds.SmoothMap.jetTopologySm_eq_gnTopologySm_of_locallyCompactSpace,
    LocallyConvex.Manifolds.CkMap.topology_eq_gnTopology,
    LocallyConvex.Manifolds.SmoothMap.topology_eq_gnTopologySm,
    LocallyConvex.Manifolds.CkMap.chartTopology_eq_gnTopology,
    LocallyConvex.Manifolds.SmoothMap.chartTopologySm_eq_gnTopologySm}
  \leanok
  \uses{thm:jet-topology-gn, def:gn-topology, def:jet-topology}
  Theorems~\ref{thm:mapping-topology-gn} and~\ref{thm:jet-topology-gn} hold for a target $N$ in
  any universe; only $M$ and the models share one universe. In particular, the topology of
  $C^\ell(M, N)$ is that of \cite[Def.~4.1.2]{GN26}, $\ell \in \N_0 \cup \{\infty\}$, for every
  Hausdorff $M$ \cite[Thm.~4.2.2(b)]{GN26}, and so is the chart topology when the model of $M$ is
  locally compact. Moreover, \cite[Def.~4.1.2]{GN26} is natural in the target: a diffeomorphism
  $\varphi \colon N \to N'$ induces a homeomorphism $f \mapsto \varphi \circ f$ between the
  topologies of \cite[Def.~4.1.2]{GN26} on $C^\ell(M, N)$ and $C^\ell(M, N')$.
\end{theorem}

\begin{proof}
  \leanok
  The iterated tangent bundles run over bundled manifolds whose carriers live in the universe
  $\max(u, v)$ of their tangent bundles, $u$ that of the models. A target $N$ in another universe
  is replaced by $\mathrm{ULift}\, N$, a manifold with the charts of $N$; $f \mapsto
  \mathrm{up} \circ f$ is a homeomorphism $C^\ell(M, N) \to C^\ell(M, \mathrm{ULift}\, N)$ for the
  jet topologies, the jets in corresponding charts being equal. Naturality follows from the chain
  rule $\T^j(\varphi \circ f) = \T^j\varphi \circ \T^j f$, with $\T^j\varphi$ a homeomorphism of
  inverse $\T^j(\varphi^{-1})$; applied to $\mathrm{up} \colon N \to \mathrm{ULift}\, N$, it shows
  that the transported topology is that of \cite[Def.~4.1.2]{GN26} whenever $N$ is itself a
  carrier. We do not claim this argument is new.
\end{proof}
